% Generado por scripts/70_tablas_figuras.R a partir de resultados-v1. No editar a mano.
\begin{table}[htbp]
\centering
\small
\caption{Tabla 6.12—Efecto según la intensidad del tratamiento}
\label{tab:tabla_6_12}
\begin{threeparttable}
\begin{tabular}{@{}lcccc@{}}
\toprule
 & (1) & (2) & (3) & (4) \\
 & Sin covariables & Con covariables & Tratados & IC 95 \% / Relativo \\
\midrule
\multicolumn{5}{@{}l}{\textit{Panel A. Tratamiento continuo: efecto promedio sobre la distribución de dosis}} \\
Dosis de bosque & 0,178 & — & 26 & [$-$0,21; 0,57] \\
 & (0,200) &  &  &  \\
Dosis de área & 0,178 & — & 26 & [$-$0,21; 0,57] \\
 & (0,200) &  &  &  \\
Dosis de bosque, en hectáreas & 185,4 & — & 26 & [$-$183; 553] \\
 & (187,8) &  &  &  \\
Dosis de bosque, controles regionales & 0,173 & — & 26 & [$-$0,23; 0,57] \\
 & (0,205) &  &  &  \\
Dosis híbrida (sensibilidad) & 0,148 & — & 30 & [$-$0,23; 0,53] \\
 & (0,195) &  &  &  \\
\addlinespace
\multicolumn{5}{@{}l}{\textit{Panel B. Gradiente exploratorio: cambio del efecto al pasar de dosis 0 a 1 (dosis de bosque)}} \\
Referencia g $-$ 1 & $-$0,61 & — & 26 & [$-$3,42; 2,19] \\
 & (1,43) &  &  & p = 0,50 \\
Referencia g $-$ 1, efectos fijos de cohorte & $-$0,49 & — & 26 & [$-$3,28; 2,30] \\
 & (1,42) &  &  & p = 0,70 \\
Referencia g $-$ 5 a g $-$ 1 & 1,19 & — & 26 & [$-$0,28; 2,66] \\
 & (0,75) &  &  & p = 0,13 \\
\addlinespace
\multicolumn{5}{@{}l}{\textit{Panel C. Terciles de la dosis de bosque, ATT balanceado e = 0--5 (ha/año)}} \\
Tercil bajo & $-$124,5 & 311,9 & 8 & $-$25,6 \% \\
 & (148,9) & (454,3) &  &  \\
Tercil medio & 294,2 & 153,4 & 9 & 60,9 \% \\
 & (317,7) & (146,5) &  &  \\
Tercil alto & 475,0 & $-$173,9 & 9 & 30,3 \% \\
 & (915,8) & (105,6) &  &  \\
\bottomrule
\end{tabular}
\begin{tablenotes}[flushleft]
\footnotesize
\item \textit{Notas:} Dosis calculada con los polígonos de deforestación evitada de Karnik et al. (2025) intersectados con los límites municipales: dosis de área, fracción del municipio dentro de la unión de los polígonos; dosis de bosque, fracción del bosque natural de 2000 (MapBiomas) dentro de ella. Solo tienen dosis los 26 tratados cubiertos por polígonos. Panel A: especificación principal del análisis de intensidad, Callaway, Goodman-Bacon y Sant’Anna (2024), lineal en la dosis, con los nunca tratados como comparación y la pérdida dividida por la media propia de 2001--2012, salvo la fila en hectáreas; el paquete (contdid 0.1.1) no admite covariables. Con una curva lineal, el efecto promedio no depende de la medida de dosis. La dosis híbrida usa el área declarada (truncada en 1) para los tratados sin polígono. En el estudio de eventos asociado, 3 de 17 coeficientes previos tienen un intervalo que excluye el cero (e = $-$16, $-$14 y $-$5). Panel B: pendiente de una regresión del efecto de cada municipio tratado (su cambio en e = 0--5 frente a la referencia, menos el de los nunca tratados en los mismos años) sobre su dosis; errores HC2 entre paréntesis y valor p por permutación de la dosis (9.999 permutaciones). Es exploratoria: no equivale al estimador del panel A y su lectura causal exige supuestos adicionales sobre las tendencias de municipios con distinta dosis. Con la dosis de área, las pendientes son $-$0,66 (1,52; IC 95 \%: [$-$3,63; 2,32]; p = 0,47) con la referencia g $-$ 1 y 1,17 (0,77; IC 95 \%: [$-$0,34; 2,68]; p = 0,14) con g $-$ 5 a g $-$ 1. Panel C: robustez del gradiente con el estimador binario; terciles calculados entre todos los tratados con dosis y comparados con los mismos nunca tratados; columna (4), ATT agregado sin covariables sobre la media previa del tercil. Con covariables, entre la mitad y tres cuartos de los pares cohorte-año de cada tercil no pueden estimarse, por lo que la columna (2) es solo indicativa. Con el área declarada, la correlación de rangos con la dosis por polígonos es $-$0,03 y solo el 32 \% de los municipios queda en el mismo tercil. Con los terciles de la dosis de área, que coinciden con los de bosque en el 87 \% de los municipios, el efecto sin covariables es de $-$124,5, 215,9 y 553,2 hectáreas en los terciles bajo, medio y alto.
\item \textit{Fuente:} Karnik et al. (2025); MapBiomas Colombia; Global Forest Change v1.12; cálculos propios.
\end{tablenotes}
\end{threeparttable}
\end{table}
